\documentclass[12pt]{article}
\usepackage[utf8]{inputenc}
\usepackage[T1]{fontenc}
\usepackage[margin=1in]{geometry}
\usepackage{setspace}
\usepackage{xcolor}
\usepackage[colorlinks=true, urlcolor=blue!60!black, linkcolor=blue!60!black]{hyperref}
\setlength{\parindent}{0em}
\setlength{\parskip}{0.6em}
\renewcommand{\baselinestretch}{1.0}
\begin{document}
\begin{center}
{\LARGE \textit{Garcetti v.\ Ceballos}, 547 U.S.\ 410 (2006)\par}
\vspace{0.5em}
{\large Case Brief --- Two-Page Edition\par}
\end{center}

\section{Memory Jogger}
A prosecutor's internal memo questioning a warrant affidavit was official-duty speech, unprotected by the First Amendment (\href{https://supreme.justia.com/cases/federal/us/547/410/}{Justia}).

\section{Detailed Case Facts}
Deputy District Attorney Richard Ceballos concluded that a search-warrant affidavit contained serious misrepresentations and wrote a memo recommending dismissal of the case. His supervisors proceeded with the prosecution; Ceballos alleged they then retaliated by reassigning him, transferring him, and denying him a promotion. He sued under \S~1983.

\section{Procedural History}
The district court granted summary judgment for the supervisors; the Ninth Circuit reversed under \textit{Pickering} and \textit{Connick}. After reargument, the Supreme Court reversed.

\section{Judicial Votes}
Reversed 5--4. Kennedy wrote for Roberts, Scalia, Thomas, and Alito. Stevens, Souter (with Stevens and Ginsburg), and Breyer dissented.

\section{Holding}
Speech made pursuant to official duties is not citizen speech and receives no First Amendment protection from employer discipline. The Court reserved whether the rule applies to scholarship or teaching.

\section{Analysis of Opinions}
\textbf{Majority:} official-duty speech is work product the employer ``commissioned or created''; protecting it would invite judicial supervision of agencies, and whistleblower laws offer other safeguards. \textbf{Stevens:} the answer should be ``Sometimes,'' not ``Never.'' \textbf{Souter:} protect official-duty speech on matters of ``unusual importance.'' \textbf{Breyer:} balance where professional and \textit{Brady} duties apply.

\section{Examples: Future Applications}
\textbf{Same-Side:} a building inspector demoted for an internal fire-code report has no claim. \textbf{Opposite-Side:} subpoenaed trial testimony about fraud learned on the job is citizen speech (\textit{Lane v.\ Franks}). \textbf{Fence-Sitter:} a public professor penalized for scholarship, the question the Court reserved.

\section{Critique}
Critics say \textit{Garcetti} punishes employees who report wrongdoing up the chain while leakers may keep protection, and leaves them to patchy whistleblower statutes. Supporters say it keeps judges from managing government workplaces.

\section{Key Quotations}
\begin{quote}
``We hold that when public employees make statements pursuant to their official duties, the employees are not speaking as citizens for First Amendment purposes, and the Constitution does not insulate their communications from employer discipline.''
\end{quote}
\noindent --- Justice Kennedy, for the Court, 547 U.S.\ at 421.

\end{document}
